% Generated draft: requires mathematical and source review.
\documentclass[11pt]{article}
\usepackage[margin=25mm]{geometry}
\usepackage{fontspec}
\setmainfont{Latin Modern Roman}
\usepackage{amsmath,amssymb,mathtools,unicode-math}
\setmathfont{Latin Modern Math}
\usepackage{enumitem,needspace,xurl,hyperref}
\setlength{\parindent}{0pt}
\setlength{\parskip}{5pt}
\setlength{\emergencystretch}{3em}
\newcommand{\N}{\mathbb{N}}
\newcommand{\Z}{\mathbb{Z}}
\newcommand{\Q}{\mathbb{Q}}
\newcommand{\R}{\mathbb{R}}
\newcommand{\C}{\mathbb{C}}
\newcommand{\F}{\mathbb{F}}
\newcommand{\A}{\mathbb{A}}
\newcommand{\PP}{\mathbb{P}}
\newcommand{\E}{\mathbb{E}}
\newcommand{\eps}{\varepsilon}
\newcommand{\dd}{\,\mathrm{d}}
\begin{document}
\section*{UnsolvedMath ID: 30006054}
\textbf{Source problem code:} catalogue entry
\subsection{Definitions and mathematical statement}
Let \(K\subset S^{3}\) be a knot. Following Definition 1 in the supplied excerpt, let \(\mathcal{D}(K)\) denote the collection of immersed discs in \(S^{3}\) with boundary \(K\) whose only singularities are ribbon and clasp singularities. For \(D\in\mathcal{D}(K)\), write \(n_{\mathrm{clasp}}(D)\) for its number of clasp intersections. The source defines the ribbon clasp number by
\[
c_r^{*}(K)=\min\bigl\{n_{\mathrm{clasp}}(D):D\in\mathcal{D}(K)\bigr\}.
\]
This reproduces the supplied definition at the level of its stated terminology. The excerpt does not specify the category of immersions, the local models and counting conventions for ribbon and clasp singularities, or the convention if the collection is empty. Its disc-immersion notation is also corrupted in the supplied transcription; the description above follows the accompanying words.

The remaining passage reports that the authors had not yet succeeded in applying their methods to examples and asks whether their Proposition 2 and Corollary 3 can settle open questions about the ribbon clasp number. Neither result nor any particular target question is included in the supplied record. Consequently, no self-contained conjecture, quantified application statement, or requested value of \(c_r^{*}(K)\) can be recovered. The application question is retained as an unresolved research direction, with the intended methods and target questions unspecified.
\subsection{Short English statement}
The ribbon clasp number counts the fewest clasp intersections among immersed discs bounded by a knot and having only ribbon and clasp singularities, subject to conventions absent from the excerpt. The source asks whether two referenced results can resolve open questions about this number, but supplies neither those results nor a specific target problem. This is a statement-only draft requiring human mathematical review.
\subsection{Sources}
\begin{itemize}
\item Supplied local problem registry: permanent UnsolvedMath ID \(30006054\), title ``Ribbon-Clasp Obstruction Methods,'' source problem code \texttt{catalogue entry}, collection ``Oberwolfach Reports Open Problems.'' No individual proposer is identified; the proposed year is recorded as 2024.
\item Supplied report citation: \emph{MATRIX-MFO Tandem Workshop: Invariants in Low-Dimensional Topology: Combinatorics, Geometry, and Computation} (2025), Oberwolfach Reports, pp.\ 2511--2550, DOI \url{https://doi.org/10.4171/owr/2024/43}. The supplied excerpt includes a page marker for p.\ 2529 and the text ``Invariants and Structures in Low-Dimensional Topology.'' These bibliographic details are retained as supplied and have not been reconciled independently.
\item Supplied literature source: ``Oberwolfach Report 43/2024,'' \url{https://doi.org/10.4171/owr/2024/43}. Literature assessment date: \texttt{2026-08-22}. This is supplied dated curation, not an independently verified literature assessment. It labels the research direction open, reports that the proposed ribbon-clasp obstruction had not yet been applied to examples, and states that no later proof of the extracted problem was located. The supplied background expressly calls for independent verification of time-sensitive status and cited claims.
\item The registry records the statement status as \texttt{unrecoverable}. Its supplied verification note says the report discusses attempted applications of proved propositions without formulating a separate self-contained open problem; no clean statement was promoted, and the original record was retained for provenance. These are supplied curatorial qualifications.
\end{itemize}
\end{document}
